\chapter{\UsamiIntroductionTitle}
% 【この章の目的】研究全体の問いと貢献を，専門外の読者にも追える形で示す．
% 内容入り例：examples/chapters/01-introduction-example.tex
% 推奨順序：背景 → 学術的課題 → 目的 → アプローチ → 貢献 → 論文構成．
% 「近年注目されている」だけで始めず，根拠となる文献・事実を示す．

\section{研究背景}
% 対象分野と問題の重要性を記載する．

\section{研究目的}
% 解く問い，対象範囲，達成条件を明示する．

\section{本研究の貢献}
% 技術的貢献，実験的知見，データ・実装上の貢献を分けて箇条書きにしてもよい．

\section{本論文の構成}
% 各章の役割を1文ずつ説明する．
